% !TEX root = ../main.tex

\chapter{结论与展望}

\section{阶段性进展}

截至中期检查，本课题围绕“下游收缩流道对横向射流火焰熄火边界影响规律”这一核心问题，从实验平台建设、测量方法建立、初步实验探索和测量手段改进四个方面开展了系统的研究工作，已完成的主要工作如下：

（1）实验平台建设与调试。完成了轴向分级燃烧器主体结构的安装与气密性检查，实现了燃烧器的稳定运行。完成了气路系统PID闭环控制的调试，解决了压缩机补压引起的压力波动问题，将一级空气流量波动控制在设定值 \(\pm 1.5\%\) 以内。基于LabVIEW平台开发了多通道RS485数据采集程序，实现了流量、温度、火焰信号的同步实时记录与存储。

（2）收缩体的设计加工。完成了五次曲线和Witoszynski曲线两种型线的计算与绘制，设计了三种收缩段，使用SolidWorks建立三维模型，采用304不锈钢材料加工完成三种收缩体，加工精度控制在 \(\pm 0.02\mathrm{mm}\)。

（3）测量手段改进探索。针对商用4-20mA探测器响应速度慢、丢失高频信息的局限性，探索了基于日盲紫外光电管C7027A的高频脉冲信号采集方案。完成了高压偏置与脉冲提取外围电路的设计，用示波器成功观测到原始脉冲信号（幅度2-3V、频率约6kHz）；编写了Python脉冲检测算法，设置阈值2.0V、死区时间5μs，实现了批量波形文件的自动脉冲计数，验证了方案的技术可行性。

（4）同温工况筛选。通过多轮筛选确定了12个同温工况（温度约 \(1069\pm 5^{\circ}C\) ），覆盖从360SLM到970SLM的完整流量范围，有效解耦了主流温度和速度的影响，为后续对比实验提供了可对比的基础条件。

（5）吹熄测量判据建立。将二级燃料从甲烷优化为乙烯，解决了中高流量工况下火焰可见度不足的问题。通过Python高频串口读取将探测器采样频率从5Hz提升至20Hz，证实了吹熄前存在信号波动过程。建立了基于占空比的四级吹熄判据体系，使火焰稳定性从主观判断转变为可重复计算的量化指标。

（6）初步对比实验。在820SLM同温工况下完成了有无收缩段的对比实验（无收缩段占空比0.14，60mm收缩段0.24，提升约 \(71\%\) ），并完成了60mm与32mm两种收缩段在五个特征工况点（360、510、650、820、970 SLM）的占空比测试，获得了两种构型的占空比随空气流量变化曲线。

\section{阶段性结论}

综合以上工作，可以得出以下阶段性结论：

（1）基于C7027A的脉冲信号采集方案能够获取商用4-20mA探测器无法提供的高频原始信号（响应速度微秒级、信息维度包含强度与频率），在原理验证层面已获成功，验证了通过自行设计外围电路获取原始紫外脉冲信号的技术可行性，为后续替代商用探测器进行高频火焰信号采集提供了可行的技术路径。

（2）同温工况筛选是确保对比实验有效性的关键前提。通过±5K容差筛选出了同温工况，使后续实验能够独立考察主流速度对熄火边界的影响，避免了温度变化对实验结果的干扰。

（3）基于占空比的四级吹熄判据能够有效区分火焰的不同稳定状态（持续燃烧、近场频闪、下游偶发复燃、完全吹熄），将火焰稳定性从主观判断转变为可重复计算的量化指标，为不同构型之间的系统对比提供了统一的评价尺度。

（4）基于现有实验数据发现，收缩段的加入对火焰稳定性有积极影响。无收缩段构型占空比为0.14，60mm五次曲线收缩段构型占空比为0.24，占空比提升约71%。在实验涉及的两种收缩高度中，32mm收缩段优于60mm收缩段占空比，初步表明存在一个最优收缩比范围。最优收缩比的可能机理是：适度的主流加速增强了剪切层混合，有利于火焰稳定；但过强的加速（60mm收缩段）可能增大火焰面的拉伸应变率，使局部Karlovitz数超过临界值，反而削弱稳定效果。

\section{创新点总结}

本课题的主要创新点可以归纳为以下四个方面：

（1）测量手段创新。针对商用4-20mA紫外火焰探测器响应速度慢、无法捕捉吹熄前高频闪烁的局限，提出了基于日盲紫外光电管C7027A的脉冲信号采集方案。设计了高压偏置与脉冲提取外围电路，用示波器成功观测到幅度2-3V、频率约6kHz的原始脉冲信号，编写了Python脉冲检测算法，验证了方案的技术可行性，为后续高频火焰信号采集提供了新的技术路径。

（2）实验方法创新。建立了同温工况筛选方法，有效解耦了主流温度和速度对熄火边界的影响，使后续对比实验能够独立考察速度这一核心变量的作用。

（3）测量方法创新。建立了基于紫外探测器信号“占空比”的量化熄火判据，将火焰稳定性从主观判断转变为可重复计算的数值指标，为热流工况下横向射流火焰吹熄的定量研究提供了统一的评价尺度。

（4）研究视角创新。在热流工况下系统对比了无收缩段、 \(60\mathrm{mm}\) 和 \(32\mathrm{mm}\) 三种构型的横向射流火焰稳定性，初步发现了最优收缩比的存在，为燃气轮机燃烧室下游过渡流道的收缩设计提供了初步的实验参考。

\section{后续工作计划}

在前期以跑通为主要目标、完成实验平台搭建、同温工况筛选、初步对比实验和占空比判据建立的基础上，后续工作将以优化为重点，按以下计划推进：

2026年5-6月：在固定二级射流参数的条件下，在同一批次内完成无收缩段、 \(60\mathrm{mm}\) 和 \(32\mathrm{mm}\) 三种构型不同型线收缩段的严格对照实验，通过多轮重复实验获取具有统计意义的数据。同时补充更多中间流量工况点（如700、750、850、900SLM等）以完善占空比曲线。基于C7027A的脉冲信号采集方案方面，搭建面包板原型电路进行验证，用数据采集卡标定最佳阈值和死区时间。

2026年6-7月：获得更可靠的完整实验数据后，进行统一处理和分析，进一步完善占空比判据，绘制不同收缩构型的熄火边界曲线。完成C7027A方案的硬件替代（LM393比较器+单片机），实现脱机实时脉冲计数。

2026年7-9月：总结收缩流道对横向射流火焰熄火边界的影响规律，将自制脉冲采集模块集成到实验台中，完成学位论文撰写。